\documentclass[a4paper,11pt]{article}

\usepackage[ngerman]{babel}
\usepackage[top=3cm,bottom=3cm,left=3cm,right=3cm]{geometry} %NICHT IN KEMIE2


%\usepackage{microtype} %schöneres typesetting  mit [stretch=10] für nur 1% anstelle von 2%
%\usepackage{lmodern} %Latin Modern FONT
%\usepackage{pifont} %schönere Font



\usepackage{amsmath}
\usepackage{nicefrac}
\usepackage{amssymb}
%\usepackage{mathtools}

\usepackage{titlesec} %Um Section, etc bearbeiten zu können
\usepackage{setspace}
\usepackage{enumitem} %enumerate
\usepackage{multicol}
\usepackage[many]{tcolorbox}
\usepackage{hyperref}
\usepackage{arydshln}

%\renewcommand{\ttdefault}{lmtt} %Schriftart wird geändert zu Latin Modern Typewriter


\hypersetup{hidelinks,
			colorlinks=true,
			linkcolor=black,
			citecolor=miudiutex,
			urlcolor=miugray2
			}

\setlength{\parindent}{0pt}
\setcounter{section}{0}


\titlespacing*{\section}{0pt}{20mm}{6mm}
\titlespacing*{\subsection}{0pt}{6mm}{3mm}
\titlespacing*{\subsubsection}{0pt}{3mm}{0mm}



%\definecolor{miudiutex}{HTML}{2f3d39}
\definecolor{miudiutex}{HTML}{1d2926}
\definecolor{miugray}{HTML}{b4cccf}
\definecolor{miugray2}{HTML}{4d7365}
\definecolor{red}{HTML}{cf1f5c}
\definecolor{green}{HTML}{00A36C}
\definecolor{delphinidin}{HTML}{315794}
\definecolor{pelargonidin}{HTML}{b33807}
\definecolor{cyanidin}{HTML}{ba0746}
\definecolor{petunidin}{HTML}{b56fed}



\raggedcolumns %wichtig für Multicolumns


%================================
% SYMBOLS
%================================

\newcommand{\RA}{$\rightarrow$~}
\newcommand{\RL}{$\rightleftharpoons$~}
\newcommand{\HARP}{\ding{226}~}
%$\xrightarrow{2}$ %Pfeil mit 2 drüber


%================================
% SHORT COMMANDS (BOXES ; ETC)
%================================


\newcommand{\notiz}[1]{\begin{notizz}{}{}\begin{center}
			{\color{miugray2!50!miudiutex}#1} \end{center}\end{notizz}}
\newcommand{\zit}[2]{\begin{zitat*}{#1}{}\small #2 \end{zitat*}}

\newcommand{\frage}[2]{\begin{qst*}{#1}{}\small #2 \end{qst*}}

\newcommand{\unten}{\end{center} \tcblower \begin{center}}


\newcommand*\circled[1]{\tikz[baseline=(char.base)]{
            \node[shape=circle,draw,inner sep=0.5pt,miudiutex] (char) {\tiny #1};}}
            
\newcommand{\miusquare}{{\color{miugray2!50}\scriptsize $\blacksquare$}~}
\newcommand{\miusquarp}{{\color{petunidin!50}\scriptsize $\blacktriangleright$}~}
\newcommand{\niu}{\item[\miusquare]}
\newcommand{\nip}{\item[\miusquarp]}

%%%%% ABSAETZE

\newcommand{\pps}{\par \vspace*{1mm}}
\newcommand{\pp}{\par \vspace*{3mm}}
\newcommand{\ppbig}{\par \vspace*{8mm}}
\newcommand{\pphuge}{\par \vspace*{10mm}}

%================================
% FRAGE BOX
%================================

\tcbuselibrary{theorems,skins,hooks}
\newtcbtheorem{qst}{Frage:}
{%
	enhanced,
	breakable,
	colback = gray!15!white,
	frame hidden,
	boxrule = 0sp,
	borderline west = {2.5pt}{0pt}{red!70},
	sharp corners,
	detach title,
	before upper = \tcbtitle\par\smallskip,
	coltitle = black,
	fonttitle = \bfseries\sffamily,
	description font = \mdseries,
	separator sign none,
	segmentation style={solid, gray!50!black},
}
{th}




%================================
% ZITAT BOX
%================================

\tcbuselibrary{theorems,skins,hooks}
\newtcbtheorem{zitat}{Zitat}
{%
	enhanced,
	breakable,
	colback = gray!15!white,
	frame hidden,
	boxrule = 0sp,
	borderline west = {2.5pt}{0pt}{black!70},
	sharp corners,
	detach title,
%	before upper = \tcbtitle\par\smallskip,
	coltitle = gray!50!black,
	fonttitle = \bfseries\sffamily,
	description font = \mdseries,
%	separator sign none,
	segmentation style={solid, gray!50!black},
}
{th}



%================================
% SIMPLE SCHWARZER RAHMEN BOX
%================================


\newcommand{\boxfr}[1]{
	\begin{tcolorbox}[
	drop shadow southeast,
	enhanced,
	colframe=black!50,
	colback=white,
	boxrule = 1pt, 
%	toprule=0pt, bottomrule=0pt,rightrule=0pt,leftrule=0pt,
	boxsep=5pt,
	arc=1pt,
	]
	\begin{center}
	#1
	\end{center}
	\end{tcolorbox}
}

\definecolor{miudiutex}{HTML}{1d2926}
\definecolor{miugray}{HTML}{b4cccf}
\definecolor{miugray2}{HTML}{4d7365}
\definecolor{white2}{HTML}{f5f7f8}
\definecolor{red}{HTML}{cf1f5c}
\definecolor{green}{HTML}{00A36C}
\definecolor{delphinidin}{HTML}{315794}
\definecolor{uniblau}{HTML}{004e9f}
\definecolor{unigelb}{HTML}{fcba00}
\definecolor{unigrau}{HTML}{909085}
\definecolor{pelargonidin}{HTML}{b33807}
\definecolor{cyanidin}{HTML}{ba0746}
\definecolor{paeonidin}{HTML}{d15ca6}
\definecolor{petunidin}{HTML}{b56fed}
\definecolor{malvidin}{HTML}{8a2d8c}

\newcommand{\metermilli}{\mskip3mu \text{mm}}
\newcommand{\cm}{\mskip3mu \text{cm}}
\newcommand{\m}{ \text{m}}
\newcommand{\lite}{ \text{l}}
\newcommand{\kg}{ \text{kg}}
\newcommand{\g}{ \text{g}}
\newcommand{\km}{ \text{km}}
\newcommand{\h}{ \text{h}}
\newcommand{\s}{ \text{s}}
\newcommand{\tonne}{ \text{t}}
\usepackage[utf8]{inputenc}
\usepackage[T1]{fontenc}
\usepackage{lmodern}

\usepackage[
    activate={true,nocompatibility},
    final,
    tracking=true,
    kerning=true,
    spacing=true,
    factor=1100,
    stretch=30,
    shrink=20
]{microtype}
\usepackage{pdfpages}
\usepackage{awesomebox}
\usepackage{marginnote}
\tcbuselibrary{documentation}
\usepackage{sidenotes}
\usepackage{import}
\usepackage{xifthen}
\usepackage{transparent}
\usepackage[table]{xcolor}


\newcommand{\abbicon}{%
  \protect\tikz[baseline=0.2mm,scale=0.8]{%
    % Das blaue Quadrat mit abgerundeten Ecken
    \fill[uniblau, rounded corners=1.2pt] (0,0) rectangle (0.8em, 0.8em);
    % Der weiße Kreis in der Mitte
    \fill[white] (0.4em, 0.4em) circle (0.12em);
  }%
} 
\usepackage[labelfont={bf},font={color=miudiutex,small},figurename={\abbicon $~$Abbildung}]{caption}


\newcommand{\bckgrnd}{
\AddToHookNext{shipout/background}{
\begin{tikzpicture}[remember picture, overlay]
 \draw[draw=none,fill=uniblau!70, shift={(current page.north west)}] (-3,0) rectangle (23,-3.75);
 \draw[draw=none,fill=miugray, shift={(current page.north east)}] (-7,0) rectangle (3,-3.75);
\end{tikzpicture}
}
}



\newcommand{\incfig}[2]{%
\def\svgwidth{#2\columnwidth}
\import{C://LaTeX-Files/pngs/}{#1.pdf_tex}
}

\renewcommand{\textbf}[1]{\textsf{{\bfseries {#1}}}}



\titleformat{\section}[hang]
		{\LARGE \color{uniblau!75!black} \sffamily \bfseries}
		{\thesection}
		{6mm}
		{}
		[]
\titleformat{\subsection}[hang]
		{\large \sffamily \bfseries \color{black}}
		{\thesubsection}
		{4mm}
		{{\color{miugray}$\blacksquare ~$}}

\titleformat{\subsubsection}[block]
		{\normalsize \bfseries \sffamily \color{miudiutex}}%
		{\normalsize \color{black} \thesubsubsection}
		{2mm}
		{{\color{miugray}$\blacksquare ~$}}
		[ \color{black}]
		
\titleformat{\paragraph}
		{\normalsize \bfseries \sffamily \color{black}}
		{\footnotesize \theparagraph}
		{1mm}
		{}

\pagestyle{empty}
\titlespacing*{\section}{0pt}{20mm}{12mm}
\titlespacing{\section}{0pt}{20mm}{12mm}

\titlespacing*{\subsection}{0pt}{14mm}{4mm}
\titlespacing{\subsection}{0pt}{14mm}{4mm}

\titlespacing*{\subsubsection}{0pt}{6mm}{3mm}
\titlespacing{\subsubsection}{0pt}{6mm}{3mm}

\titlespacing*{\paragraph}{0pt}{6mm}{2mm}
\titlespacing{\paragraph}{0pt}{6mm}{2mm}


\let\oldmarginfigure\marginfigure
\let\endoldmarginfigure\endmarginfigure

\let\oldmarginfigure\marginfigure
\let\endoldmarginfigure\endmarginfigure

\renewenvironment{marginfigure}[1][0pt] % [1] steht für 1 Argument, [0pt] ist der Standardwert
  {\oldmarginfigure[#1]\captionsetup{labelfont={sf,bf,color=uniblau!75!black},font={color=miudiutex,small}}}
  {\endoldmarginfigure}
  
  
  
  
\renewcommand{\labelitemi}{\color{uniblau!70}$\bullet$}
\renewcommand{\labelitemii}{\color{uniblau!70}$\cdot$}
\newcommand{\plmtsquare}{{\color{uniblau!70}$\blacksquare~$}}
\newcommand{\splmt}[1]{{\color{uniblau!75!black} \sffamily #1}}
\newcommand{\fplmt}[1]{{\color{uniblau!75!black} \sffamily \bfseries #1}}

\newcommand{\antwort}[1]{
\begin{tcolorbox}[
	enhanced,
	enforce breakable,
	standard jigsaw,
	boxrule=0.7pt,
	colframe=black,
	sharp corners,
	boxsep=5pt,
	title=\bfseries \sffamily Antwort,
	opacityback=0,
	coltitle=miudiutex,
	attach boxed title to top text left={yshift=-\tcboxedtitleheight/2},
	boxed title style={colback=white,colframe=white},
	bottomrule at break=0pt,
	toprule at break=0pt,
	pad at break=16mm,
]
	\begin{center}
	#1
	\end{center}
	\end{tcolorbox}
}




\rowcolors{2}{miugray!50}{miugray!50}
\arrayrulecolor{white} % Weiße Gitterlinien
\setlength{\arrayrulewidth}{1.5pt} % Etwas dickere Linien
\renewcommand{\arraystretch}{1.4} % Mehr Zeilenhöhe für bessere Lesbarkeit}


%\setlist[itemize]{itemsep=0mm}
\setlist[enumerate]{itemsep=0mm}








\begin{document}
\newgeometry{top=1.8cm,bottom=4cm,left=1.5cm,right=7.5cm,marginparsep=-2.00cm,marginparwidth=4.75cm}


\bckgrnd


\begin{center}
\LARGE
\color{white}
\textbf{Übung 01} \par \vspace*{2mm}
\large \color{miugray} \textbf{Prozessbezogene Lebensmitteltechnologie}
\end{center}
\par 
\vspace*{6mm}


\color{miudiutex}
\section*{Strömungslehre I}

\subsection*{Aufgabe 1}
Ein Fluid 
%\splmt{Fluid} 
%\sidenote{Fluide können sowohl gasförmig als auch flüssig sein} 
durchströmt ein Rohr mit drei unterschiedlichen Querschnittsflächen. 
Im ersten Bereich beträgt die Querschnittsfläche 20 cm$^2$ und die Fließgeschwindigkeit des Fluids beträgt 3 cm/s.
%\marginnote[]{ \small
%%\color{uniblau} 
%{\color{uniblau!75!black} \bfseries \sffamily Zusatzinfo:} \par 
%Fluide können sowohl gasförmig als auch flüssig sein \par 
%$A_1 \cdot v_1 = A_2 \cdot v_2$
%}
\begin{itemize}
\item Im zweiten Bereich beträgt die Querschnittsfläche 10 cm$^2$. Wie hoch ist die Geschwindigkeit des Fluids hier?
\item Im dritten Bereich beträgt die Querschnittsfläche 25 cm$^2$. Wie hoch ist die Geschwindigkeit des Fluids hier?

\end{itemize}

 

Kontinuitätsgleichung
\begin{align*}
A_1 \cdot v_1 = A_2 \cdot v_2 \\
v_2 = \frac{A_1 \cdot v_1}{A_2}
\end{align*}
Einsetzen a) und b):
\begin{align*}
a) &&v_2 &= \frac{20\cm^2 \cdot 3\mskip3mu \cm/\text{s}}{10\mskip3mu \cm^2} &&= 6\cm/\text{s} \\
b) &&v_2 &= \frac{20\cm^2 \cdot 3\mskip3mu \cm/\text{s}}{25\mskip3mu \cm^2} &&=2,4\cm/\text{s}
\end{align*}

\pagebreak
 
 





%\pagebreak
%\newgeometry{top=3cm,bottom=3cm,left=4cm,right=4cm}
\subsection*{Aufgabe 2}
 
\begin{marginfigure}[+2cm]
 
\begin{marginfigure}[-2cm]
 
\incfig{venturi}{.9}
\caption{Beim Venturi-Rohr lässt sich die Fließgeschwindigkeit über den Höhenunterschied der angesaugten Flüssigkeit in einem Rohr ablesen, welches die Verengung mit dem normalen Rohrdurchmesser verbindet.}
\end{marginfigure}

In einer Wasserstrahl-Vakuumpumpe, die wie ein \splmt{Venturi-Rohr} aufgebaut ist, fließen 4 L/min Wasser. Der Durchmesser der Zuleitung (und Ableitung) beträgt \mbox{14 mm}, der Durchmesser an der engsten Stelle 3 mm. Der absolute Druck am Ausgang der Wasserstrahl-Vakuumpumpe entspricht einem Umgebungsdruck von 1 bar. Der geodätische Höhenunterschied kann vernachlässigt werden. 
\begin{itemize}
\item Welcher Unterdruck wird an der engsten Stelle erreicht? 
\item Welcher absolute Druck herrscht an der engsten Stelle?
\end{itemize}

 

\begin{align*}
\dot{V} &= ~\text{const.} \\
\dot{V} &= v \cdot A \qquad \rightarrow v = \frac{\dot{V}}{A} \\ \\
v_1 &= \frac{\dot{V}}{A_1} \\
	&= \frac{4~\text{L/min}}{\pi \cdot (0,007~\text{m})^2} \\
	&= \frac{6,\overline{6} \cdot 10^{-5} ~\text{m$^3$}}{1,54 \cdot 10^{-4}~\text{m}^2 \cdot \text{s}} \\
v_1 &= 0,43\mskip3mu \m/\text{s} \\ \\
v_2 &= \frac{v_1 \cdot A_1}{A_2}\\
&= v_1 \cdot \frac{d_1^2}{d_2^2} \\
&= 0,43\mskip3mu \text{m/s} \cdot \frac{14^2}{3^2} \\
&=9,36\mskip3mu \text{m/s}
\end{align*}

Bernoulli:
\begin{equation}
p_1 + \frac{\rho \cdot v_1^2}{2} = p_2 + \frac{\rho \cdot v_2^2}{2}
\end{equation}
\begin{align*}
p_1 - p_2 = \frac{\rho \cdot v_2^2}{2} -  \frac{\rho \cdot v_1^2}{2} 
\end{align*}

\pp



Werte einsetzen: 
\begin{align*}
p_1 - p_2 &= \frac{1000 \mskip3mu \text{kg/m}^3 \cdot (9,36\mskip3mu \text{m/s} )^2}{2} -  \frac{1000 \mskip3mu \text{kg/m}^3 \cdot (0,43 \mskip3mu \text{m/s} )^2}{2} \\
&= 43.712 \mskip3mu \frac{\text{kg}}{\m \cdot \s^2} \\
&= 43,7 \mskip3mu \text{kPa} \\ \\
p_2 &= p_1 - \Delta p \\
	&= (100 - 43,7)\text{kPa}\\ &= 56,3\mskip3mu \text{kPa}
\end{align*}

Wobei $p_2$ der Druck an der engsten Stelle und $p_1$ der Druck am Ende des Rohres ist.\par
Absoluter Druck $p_2$ = 56,3 \text{kPa} \par 
Unterdruck $\Delta p$ = 43,7 \text{kPa} \par 

\pagebreak
 
 


\subsection*{Aufgabe 3}
Wie viel Zeit braucht die Entleerung eines Rührkessels unter der Wirkung der \mbox{Erdschwere}, wenn der Bodenablass die \splmt{NW} 80 mm aufweist? Der Rührkessel ist zu Beginn mit \mbox{10.000 L} Wasser gefüllt. Der Innendurchmesser des Rührkessels beträgt 2,4 m. Der gewölbte Boden des Klöppers ist bei der Berechnung zu vernachlässigen, ebenso der Druckverlust beim Einströmen in die Rohrleitung.




\marginnote[]{ \small Die {\color{uniblau!75!black} \bfseries \sffamily Nennweite (NW)} beschreibt den ungefähren Innendurchmesser eines Rohres und wurde eingeführt zur Vereinheitlichung von Bauteilen nach einer ISO-Norm}[-2.5cm]

 


\begin{align*}
&A_1 &&= \pi \cdot (\frac{2,4\mskip3mu \m}{2})^2 &&= 4,52\mskip3mu \m^2 \\
&A_2 &&= \pi \cdot (\frac{0,08\mskip3mu \m}{2})^2 &&= 0,005\mskip3mu \m^2 \\
&h_a &&= \frac{V}{A_1} = \frac{10\mskip3mu \m^3}{4,52~\text{m}^2} &&= 2,21\mskip3mu \m^2 
\end{align*}

\ppbig
Torricelli:
\begin{equation}
t = \sqrt{\frac{2}{g}} \cdot \frac{A_1}{A_2} \cdot (\sqrt{h_a} - \sqrt{h_w})
\end{equation}

Einsetzen:
\begin{align*}
t &= \sqrt{\frac{2}{g}} \cdot \frac{4,52\mskip3mu \m^2}{0,005\mskip3mu \m^2} \cdot \sqrt{2,21\mskip3mu \m^2} \\
&= 604 ~\text{s}
\end{align*}




\pagebreak
 
 
%\newgeometry{top=3cm,bottom=3cm,left=4cm,right=4cm}
\pagebreak

\subsection*{Aufgabe 4}
Der Anlieferer einer neuen Laborausrüstung für die Uni Bonn hat bei seiner viel zu späten Ankunft am Abend in seiner Eile zwei voll befüllte Weinfässer am Boden beschädigt. Das Leck des ersten Fasses hat eine Querschnittsfläche von 0,03 cm$^2$, das zweite Fass eine Querschnittsfläche von 0,05 cm$^2$. Beide Fässer haben einen Durchmesser von 50 cm und fassen ein Füllvolumen von 200 L. \par 
%\textit{Gehen Sie von einem linearen Verlauf der Entleerung aus.}

\begin{marginfigure}[0cm]
%\begin{figure}[ht]
\centering
\incfig{fass}{.7}
%\caption{Ein Weinfass der Uni Bonn hat ein Leck und läuft aus.}
%\end{figure}
\end{marginfigure}
 
\vspace*{0mm}
 
  
\begin{itemize}
\item Die beschädigten Fässer wurden erst am nächsten Morgen, 8 Stunden nach der Anlieferung, von der Laborhilfskraft entdeckt und sollten heute eigentlich auf ihren Anthocyangehalt untersucht werden. Haben sich beide Fässer über Nacht entleert?
 

Füllstand $h$:
\begin{align*}
h = \frac{V}{\pi \cdot r^2} = \frac{0,2~\text{m}^3}{\pi \cdot (0,25~\text{m})^2} = 1,0186~\text{m}
\end{align*}

\begin{align*}
t &= \sqrt{\frac{2}{g}} \cdot \frac{A_1}{A_2} \cdot \left(\sqrt{h_{\alpha}} - \sqrt{h_w} \right) \\
 &= \sqrt{\frac{2}{9,81~\text{m/s}^2}} \cdot \frac{1.963,5~\text{cm}^2}{0,03~\text{cm}^2} \cdot \left(\sqrt{1,0186~\text{m}} - 0 \right) \\
 &= 29.825,79~\text{s} \qquad \approx 8,28~\text{h}
\end{align*}
Das Fass mit dem kleinen Leck ist nach 8 Stunden noch nicht komplett entleert.
\begin{align*}
t &= \sqrt{\frac{2}{9,81~\text{m/s}^2}} \cdot \frac{1.963,5~\text{cm}^2}{0,05~\text{cm}^2} \cdot \left(\sqrt{1,0186~\text{m}} - 0 \right) \\
 &= 17.895,48~\text{s} \qquad \approx 4,97~\text{h}
\end{align*}
Das Fass mit dem größeren Leck ist bereits nach guten 5 Stunden komplett entleert.

  
\item Wie verhält sich die Entleerungszeit bei Änderung der Querschnittsfläche des Ausflusses im Gegensatz zu einer der Änderung der Anfangsfüllhöhe?


 
Die Entleerungszeit verhält sich umgekehrt proportional zur Querschnittsfläche des Ausflusses, während sie proportional zur Füllhöhe der Wurzel ist. \pp 
Eine Verdopplung der Entleerungszeit bedeutet eine \textbf{Vervierfachung} der Füllhöhe. \par 
Eine Verdopplung der Entleerungszeit bedeutet eine \textbf{Halbierung} der Ausflussfläche. 


\pagebreak
 
  

\end{itemize}






\subsection*{Aufgabe 5}
Berechnen Sie die Reynoldszahl für Wasser bei 0°C. Ist die Strömung jeweils laminar oder turbulent?

$\rho_{H_2O}$ (0°C) = 999,8 kg/m$^3$ ; $\eta$ = 1.790 $\cdot ~ 10^{-6}$ Pa $\cdot$ s. 

\begin{itemize}
       \item Geschwindigkeit 5 cm/s, Rohrinnendurchmesser 30 cm
       \item Geschwindigkeit 1 cm/s, Rohrinnendurchmesser 30 cm
       \item Geschwindigkeit 5 cm/s, Rohrinnendurchmesser 10 cm
\end{itemize}



%\tcbdocmarginnote[%
%	enhanced,
%	boxrule=0.7pt,
%	colframe=miugray,
%	sharp corners,
%	boxsep=5pt,
%	colback=miugray,
%	coltitle=white,
%	]
%{{ \bfseries \sffamily Schon gewusst?} \par 
%Wasser ist nass.
%}



 


Kritische Reynoldszahl für turbulente oder laminare Strömung:
\begin{equation}
\text{Re}_{\text{krit}} = \frac{\rho \cdot d \cdot v}{\eta} = 2.300
\end{equation}


Einsetzen:
\begin{align*}
&\text{a)} && \frac{1.000\mskip3mu \kg/\m^3 \cdot 0,3\mskip3mu \m \cdot 0,05\mskip3mu \m/\s}{1.790 \cdot 10^{-6}\mskip3mu \text{P} \cdot \s} &&\approx 8.379,9 \\
&\text{b)} && \frac{1.000\mskip3mu \kg/\m^3 \cdot 0,3\mskip3mu \m \cdot 0,01\mskip3mu \m/\s}{1.790 \cdot 10^{-6}\mskip3mu \text{P} \cdot \s} &&\approx 1.676 \\
&\text{c)} && \frac{1.000\mskip3mu \kg/\m^3 \cdot 0,1\mskip3mu \m \cdot 0,05\mskip3mu \m/\s}{1.790 \cdot 10^{-6}\mskip3mu \text{P} \cdot \s} &&\approx 2.793,3 \\
\end{align*}

Oberhalb der kritischen Reynolds-Zahl 2.300 erfolgt der Übergang zu turbulenten Strömungen.

\pagebreak
 
 


\subsection*{Aufgabe 6}
Ein Landwirt transportiert mit seinem Traktor Nitrat-Dünger zu einem nahegelegenen Feld. Beim Überfahren eines Schlaglochs löst sich die Verschlusskappe (Durchmesser = 5 cm) eines der zylinderförmigen Kanister. Dem Landwirt fällt das Abfallen der Verschlusskappe sofort auf. Er hält an und befestigt die Verschlusskappe an der Öffnung, die sich auf Höhe des Kanisterbodens befindet. Für diesen Prozess benötigt er 2 Minuten. Die verbleibende Füllhöhe schätzt der Landwirt auf circa 20 cm. Leider weiß er nicht mehr, wie hoch die Kanister anfänglich befüllt waren.\par 
Wie viele Liter Nitrat-Dünger sind bis dahin ausgelaufen, wenn der Kanister einen Durchmesser von 1 m hat? 


	
 

Zunächst ursprünglichen Füllstand des Kanisters berechnen:
\begin{align*}	
t &= \sqrt{\frac{2}{g}} \cdot \frac{A_1}{A_2} \cdot \left(\sqrt{h_{\alpha}} - \sqrt{h_w} \right) \\
t &= \sqrt{\frac{2}{g}} \cdot \frac{A_1}{A_2} \cdot \sqrt{h_{\alpha}} - \sqrt{\frac{2}{g}} \cdot \frac{A_1}{A_2} \cdot  \sqrt{h_w} \\ \\
\sqrt{h_{\alpha}} &= \frac{t + \sqrt{\frac{2}{g}} \cdot \frac{A_1}{A_2} \cdot  \sqrt{h_w}}{\sqrt{\frac{2}{g}} \cdot \frac{A_1}{A_2}}\\
h_{\alpha} &= \left( \frac{t + \sqrt{\frac{2}{g}} \cdot \frac{A_1}{A_2} \cdot  \sqrt{h_w}}{\sqrt{\frac{2}{g}} \cdot \frac{A_1}{A_2}} \right)^2 \\
h_{\alpha} &= \left( \frac{120~\text{s} + \sqrt{\frac{2}{9,81~\text{m/s}^2}} \cdot \frac{(1~\text{m})^2}{(0,05~\text{m})^2} \cdot  \sqrt{0,2~\text{m}}}{\sqrt{\frac{2}{9,81~\text{m/s}^2}} \cdot \frac{(1~\text{m})^2}{(0,05~\text{m})^2}} \right)^2 \\
&\approx 1,236~\text{m}
\end{align*}
\pp

Über den Füllstand das Volumen berechnen:
\begin{align*}
V &= \pi \cdot (0,5~\text{m})^2 \cdot (1,236 - 0,2~\text{m}) \\
&= 0,814~\text{m}^3 \\
&= 814 ~\text{dm}^3 \\
&= 814~\text{L}
\end{align*}

\pagebreak
 
 



\end{document}
